\section*{الفصل \ref{ch:g3:levers-and-scales} --- الروافع وموازين الكفّتين}

\begin{solution}{exo:g3:levers-and-scales:1}
الأرجوحة: اللوح هو \omterm{def:g3:levers-and-scales:lever}{الرافعة}، والسند الأوسط \omterm{def:g3:levers-and-scales:lever}{محور الارتكاز}.
والعتلة: القضيب هو \omterm{def:g3:levers-and-scales:lever}{الرافعة}، والجذع \omterm{def:g3:levers-and-scales:lever}{محور الارتكاز}. والمقصّ: كل
نصل مع مقبضه \omterm{def:g3:levers-and-scales:lever}{رافعة}، وكلاهما يرتكز على البرغي المركزي.
\end{solution}

\begin{solution}{exo:g3:levers-and-scales:2}
بعيدًا عن \omterm{def:g3:levers-and-scales:lever}{محور الارتكاز}: فكلما بَعُدت الدفعة كبر جداؤها ---
ودفعة صغيرة بعيدة تغلب صخرة ثقيلة قريبة.
\end{solution}

\begin{solution}{exo:g3:levers-and-scales:3}
$3 \times 4 = 12$؛ \ $2 \times 5 = 10$؛ \ $1 \times 6 = 6$. ولا
يتساوى شيء من هذه الجداءات، فلا يوازن أيّ اثنين منها الآخر.
\end{solution}

\begin{solution}{exo:g3:levers-and-scales:4}
جداء اليسار $1 \times 8 = 8$. والقطعتان تحتاجان إلى
$8 = 2 \times 4$: أي العلامة $4$. والقطع الأربع تحتاج إلى
$8 = 4 \times 2$: أي العلامة $2$.
\end{solution}

\begin{solution}{exo:g3:levers-and-scales:5}
اليسار: $3 \times 3 = 9$. واليمين: $2 \times 4 = 8$. فجداء
الجانب الأيسر أكبر، فينزل الجانب الأيسر.
\end{solution}

\begin{solution}{exo:g3:levers-and-scales:6}
مع تساوي الذراعين يعني استواء الكفّتين تساوي الوزنين ---
فالمقارنة عادلة. أما ذراع أيسر أطول فيضرب الحِمل الأيسر في مسافة
أكبر، فيكذب الميزان لصالح ذلك الجانب في كل وزن.
\end{solution}

\begin{solution}{exo:g3:levers-and-scales:7}
الكيس يزن تمامًا قدر القطع الستّ. وبنزع قطعة واحدة تصير كفّة
الكيس أثقل، فيميل الميزان نحو الكيس.
\end{solution}

\begin{solution}{exo:g3:levers-and-scales:8}
جداء توم: $2 \times 3 = 6$. وتحتاج أخته إلى $1 \times 6 = 6$:
أي أن تجلس عند العلامة $6$.
\end{solution}

\begin{solution}{exo:g3:levers-and-scales:9}
\omterm{def:g3:levers-and-scales:lever}{محور الارتكاز} هو الموضع الذي يستند فيه رأس المطرقة إلى الخشب.
ويدك تدفع عند طرف المقبض البعيد --- أي مسافة طويلة، فجداء كبير
--- بينما يقاوم المسمار قريبًا جدًّا من \omterm{def:g3:levers-and-scales:lever}{محور الارتكاز}: \omterm{def:g3:levers-and-scales:lever}{فالرافعة}
تحوّل دفعتك المريحة إلى سحبة عملاق.
\end{solution}

\begin{solution}{exo:g3:levers-and-scales:10}
التوازن يقتضي $4 \times (\text{علامة الجدّ}) = 1 \times
(\text{علامة زوي})$. وطريقتان: الجدّ عند $1$ وزوي عند $4$
($4 = 4$)، أو الجدّ عند $2$ وزوي عند $8$ ($8 = 8$).
\end{solution}

\begin{solution}{exo:g3:levers-and-scales:11}
أقلّ. فعلى الذراع الطويل يصنع الحِمل الخفيف جداءً كبيرًا، فتستوي
الكفّتان والسلعة تزن \emph{أقلّ} من القطع الأمينة على الذراع
القصير. وفحص الميزان فارغًا يفضحه: فبلا شيء في الكفّتين يميل
ميزانه أصلًا نحو الذراع الطويل.
\end{solution}
